\oldpage[II]{53}
à dire que $\mathcal{S}_y$
est une $\mathcal{O}_y$-algèbre graduée essentiellement réduite pour tout
$y\in Y$. Pour toute $\mathcal{O}_Y$-Algèbre graduée $\mathcal{S}$,
$\mathcal{S}/\mathcal{N}_+$ est essentiellement réduite.

Nous dirons que $\mathcal{S}$ est \emph{intègre} si $\mathcal{S}_y$ est un
anneau intègre pour tout $y\in Y$ et si en outre
$(\mathcal{S}_y)_+=(\mathcal{S}_+)_y\neq0$ pour tout $y\in Y$.
\end{env}

\begin{proposition}[3.1.13]
\label{II.3.1.13-fr}
Soit $\mathcal{S}$ une $\mathcal{O}_Y$-Algèbre graduée à degrés positifs.
Si $X=\operatorname{Proj}(\mathcal{S})$, le $Y$-schéma $X_{\mathrm{red}}$ est
canoniquement isomorphe à $\operatorname{Proj}(\mathcal{S}/\mathcal{N}_+)$~;
en particulier, si $\mathcal{S}$ est essentiellement réduite, $X$ est réduit.
\end{proposition}

\begin{proof}
Le fait que $X'=\operatorname{Proj}(\mathcal{S}/\mathcal{N}_+)$ est réduit
découle immédiatement de (\hyperref[II.2.4.4-fr]{2.4.4, (i)}), la propriété
étant locale~; en outre, pour tout ouvert affine $U\subset Y$,
${\varphi'}^{-1}(U)$ est égal à $(\varphi^{-1}(U))_{\mathrm{red}}$ (en
désignant par $\varphi$ et $\varphi'$ les morphismes structuraux $X\to Y$,
$X'\to Y$)~; on vérifie aussitôt que les $U$-morphismes canoniques
${\varphi'}^{-1}(U)\to\varphi^{-1}(U)$ sont compatibles avec les opérations
de restriction et définissent par suite une immersion fermée $X'\to X$ qui
est un homéomorphisme des espaces sous-jacents~; d'où la conclusion
(\hyperref[I.5.1.2-fr]{\textbf{I}, 5.1.2}).
\end{proof}

\begin{proposition}[3.1.14]
\label{II.3.1.14-fr}
Soient $Y$ un préschéma intègre, $\mathcal{S}$ une
$\mathcal{O}_Y$-Algèbre graduée quasi-cohérente telle que
$\mathcal{S}_0=\mathcal{O}_Y$.
\begin{enumerate}
  \item[\rm(i)] Si $\mathcal{S}$ est intègre
    \textup{(\hyperref[II.3.1.12-fr]{3.1.12})}, alors
    $X=\operatorname{Proj}(\mathcal{S})$ est intègre et le morphisme
    structural $\varphi:X\to Y$ est dominant.
  \item[\rm(ii)] Supposons en outre que $\mathcal{S}$ soit essentiellement
    réduite. Alors, inversement, si $X$ est intègre et si $\varphi$ est
    dominant, $\mathcal{S}$ est intègre.
\end{enumerate}
\end{proposition}

\begin{proof}
\begin{enumerate}
  \item[\rm(i)] Si $(U_\alpha)$ est une base de $Y$ formée d'ouverts affines
    non vides, il suffit de démontrer la proposition lorsque $Y$ est remplacé
    par un des $U_\alpha$ et $\mathcal{S}$ par $\mathcal{S}|U_\alpha$~: en
    effet, il en résultera d'une part que les espaces sous-jacents
    $\varphi^{-1}(U_\alpha)$ sont des ouverts irréductibles (donc non vides)
    de $X$ tels que
    $\varphi^{-1}(U_\alpha)\cap\varphi^{-1}(U_\beta)\neq\varnothing$ pour tout
    couple d'indices (puisque $U_\alpha\cap U_\beta$ contient un $U_\gamma$),
    donc que $X$ est irréductible (\hyperref[0.2.1.4-fr]{\textbf{0}, 2.1.4})~; d'autre
    part $X$ sera réduit, puisque c'est là une propriété locale, donc $X$
    sera bien intègre et $\varphi(X)$ dense dans $Y$.

    Supposons donc que $Y=\operatorname{Spec}(A)$ où $A$ est intègre
    (\hyperref[I.5.1.4-fr]{\textbf{I}, 5.1.4}) et $\mathcal{S}=\widetilde{S}$, où $S$
    est une $A$-algèbre graduée~; l'hypothèse est que pour tout $y\in Y$,
    $\widetilde{S}_y=S_y$ est un anneau gradué intègre tel que
    $(S_y)_+\neq0$. Il suffit de prouver que $S$ est un anneau \emph{intègre},
    car alors on aura $S_+\neq0$ et on pourra appliquer
    (\hyperref[II.2.4.4-fr]{2.4.4, (ii)}). Or, soient $f$, $g$ deux éléments
    $\neq0$ de $S$ et supposons que $fg=0$~; pour tout $y\in Y$ on aura donc
    $(f/1)(g/1)=0$ dans $S_y$, donc $f/1=0$ ou $g/1=0$ par hypothèse.
    Supposons par exemple $f/1=0$ dans $S_y$~; cela signifie qu'il existe
    $a\in A$ tel que $a\notin\mathfrak{j}_y$ et $af=0$~; on a alors pour \emph{tout}
    $z\in Y$, $(a/1)(f/1)=0$ dans l'anneau \emph{intègre} $S_z$, et comme
    $a/1\neq0$ (puisque $A$ est intègre), $f/1=0$, ce qui implique $f=0$.
  \item[\rm(ii)] La question étant locale sur $Y$, on peut encore supposer
    que $Y=\operatorname{Spec}(A)$, $A$ étant intègre, et
    $\mathcal{S}=\widetilde{S}$. Par hypothèse, pour tout $y\in Y$, $(S_y)_+$
    ne contient pas d'élément nilpotent $\neq0$, et il en est de même de
    $(S_0)_y=A_y$ par hypothèse~; donc $S_y$ est un anneau réduit pour tout
    $y\in Y$, et on en conclut d'abord que $S$ lui-même est réduit
    (\hyperref[I.5.1.1-fr]{\textbf{I}, 5.1.1}). L'hypothèse que $X$ est intègre
    entraîne d'autre part que $S$ est essentiellement intègre
    (\hyperref[II.2.4.4-fr]{2.4.4, (ii)}), et tout revient alors à voir que
    l'annulateur $\mathfrak{J}$ de $S_+$ dans $A=S_0$ est réduit
